\documentclass[11pt,a4paper,oneside]{article}
\usepackage[utf8]{inputenc}
\usepackage[T1]{fontenc}
\usepackage{amsmath}
\usepackage{amsfonts}
\usepackage{amssymb}
\usepackage{graphicx}
\usepackage{booktabs}
\usepackage{hyperref}
\usepackage{geometry}
\geometry{margin=1in}

\title{Short Communication: Analysis of German Gas Storage Drawdown and Supply Risk Assessment (Q1 2026)}
\author{Collaborative Agentic Core \& Abdelrhim \\ \textit{Energy Data Modeling Team}}
\date{March 15, 2026}

\begin{document}

\maketitle

\begin{abstract}
This report details the implementation and results of a hybrid predictive model developed to monitor German gas storage levels during the Q1 2026 heating season. By integrating ERA5-Land climate anomalies with industrial cluster-weighted HDD, the model identifies significant storage drawdown trends and supply-side anomalies. Our findings indicate a February 2026 Crisis Delta of 0.1789\% and highlight gas price volatility as a primary driver of storage residuals.
\end{abstract}

\section{Introduction}
Monitoring gas storage levels is critical for ensuring energy security in Germany. This project utilizes an agentic workflow to develop a high-precision predictive engine capable of detecting deviations from expected drawdown patterns, termed 'Crisis Deltas'.

\section{Methodology}
The analysis pipeline integrates three primary data streams:
\begin{enumerate}
    \item \textbf{Climate Data}: Extraction of 2m Temperature and Total Precipitation from ERA5-Land archives for the German coordinates.
    \item \textbf{Industrial Weighting}: Calculation of Heating Degree Days (HDD) weighted by the density of industrial clusters in NRW, Bavaria, and Baden-Württemberg.
    \item \textbf{Model Architecture}: A temporal-lagged XGBoost model (RMSE: 0.0045) was deployed to capture the delayed impact of market price signals on storage levels.
\end{enumerate}

\section{Results}
The model achieved a high tracking accuracy across the evaluation period. Key metrics are summarized in Table \ref{tab:results}.

\begin{table}[h]
    \centering
    \caption{Key Model Metrics (Q1 2026)}
    \begin{tabular}{ll}
        \toprule
        Metric & Value \\
        \midrule
        Root Mean Square Error (RMSE) & 0.0045 \\
        February 2026 Crisis Delta & 0.1789 \% \\
        Max Anomaly Magnitude & 5.2$\sigma$ \\
        \bottomrule
    \end{tabular}
    \label{tab:results}
\end{table}

Anomalous behavior was detected on \textbf{January 4th, 2026}, where the residual exceeded the 5$\sigma$ threshold, suggesting a non-climatic supply disruption or extreme market reaction.

\section{Conclusion}
The integration of climate anomalies with industrial demand weighting provides a robust framework for energy tracking. Future work will focus on integrating real-time market sentiment analysis to further reduce residuals.

\end{document}
